\chapter[Ueber Laienälteste]{Ueber Laienälteste oder Gemeindevorsteher}
\quellseite{22}

Nach und nach ist in der lutherischen Kirche der Gedanke fast herrschend geworden, das Institut der s.\,g. Laienältesten
(der presbyteri laici, seniores plebis) sei ein Schößling der Reformirten Kirche und der Lehre derselben von
Amt und Kirchenverfassung, hingegen dem Geiste und der Lehre der evangelisch-lutherischen Kirche fremd und
zuwider. Diesem Wahn hat schon vor mehreren Jahren Dr. Rudelbach mit ebenso großem Rechte, als Entschiedenheit
widersprochen. Dieser große Geschichtsforscher schreibt nämlich u.\,A. in einer Abhandlung über „Staatskirchenthum
und Religionsfreiheit“ (erschienen im Jahre 1850 in der Rudelbach-Guericke’schen Zeitschrift): „Es war Luther
wie Melanchthon durchaus einleuchtend, daß die synodale Verfassung mit den freien Gemeinde-Organen, welche sie
voraussetzt, am liebsten einen permanenten Regierungs-Organismus zur Seite (welcher die wesentliche Bedeutung der
Bischöfe), sowohl die zweckdienlichste, als die eigenthümliche Verfassung der evangelischen (luth.) Kirche sein
müßte. Dieses müssen wir schon hier scharf betonen, weil in der letzten Zeit eine Betrachtung sich hervorgebracht
hat, als ob das presbyteriale Element schlechterdings nicht lutherisch, sondern eigenthümlich, wesentlich Reformirt
sei — eine Behauptung, die ebenso durch die eignen klaren Zeugnisse der Reformatoren als durch die Natur der
Sache widerlegt wird; denn unmöglich kann sich eine organische Synodalverfassung bilden, ohne daß Presbyterien
ihr zur Seite stehen. ... Es ist ferner offenkundig, daß man (von Seiten Luther’s und Melanchthon’s), um die
Interessen der Kirche als eines Gesammtkörpers wahrzunehmen, Synoden nicht bloß wünschte, sondern laut forderte,
und daß man darunter durchaus nicht (wie von Stahl u.\,a. in der letzten Zeit fälschlich behauptet worden ist)
bloße Versammlungen von Geistlichen
\quellseite{23} Verstand, sondern eine wirkliche Kirchenrepräsentation;\footnote{„Siehe Melanchthons Werke III, 472.: ‚Aber in Streitigkeiten über die Lehre sollen nicht allein die Bischöfe judiciren, sondern es müssen geschickte Leute aus dem ganzen Körper der Kirche ausgewählt werden.‘“ Überhaupt behaupteten die Reformatoren, daß eine jede Tyrannei dem Wesen der Kirche widerstrebe, und daß man selbst den Schein davon vermeiden müsse. „Tyrannei“, sagt Melanchthon, „ist die Feindin der Kirche.“ (S. seinen Brief an die Nürnbergischen Prediger vom 17. Febr. 1540. Opp. ed. Bretschneider, T. III. 965. seqq.: „Auch der Bann muß wiederhergestellt werden, nicht, wie vormals, in Händeln weltlicher Dinge, sondern über offenbare Verbrechen, indem man zu diesem Gericht die Senioren in jeder Gemeinde dazu nimmt.“ Es ist dies übrigens der Punkt, wo die reine Idee der Presbyterien und die praktische Aufgabe, die man schon damals in der Errichtung der Consistorien vor Augen hatte, in den Umrissen gleichsam verschmolzen.“ Rudelbach.} wodurch die Gemeinde nicht bloß das zweideutige Recht, den Verhandlungen beizustimmen, sondern ein wahrhaft decisives Votum hätte, welches später auf die Laien-Mitglieder der Consistorien übertragen ward.“ (S. 396–8.)

Je wichtiger nun gerade hier in Amerika, wo Kirche und Staat getrennt ist und die kirchlichen Gemeinden unabhängig und selbstständig neben den bürgerlichen dastehen, nothwendig ein Institut sein muß, wie das sogenannte Laienältesten, und je mehr man in der rechten Anschauung desselben bald zur Rechten bald zur Linken weit abweicht: um so erwünschter dürfte es manchem hier sein, über dieses Institut, dessen biblische Begründung, Stellung, Wirkungskreis und Gewalt, die Stimme unserer rechtgläubigen Lehrväter zu vernehmen. Wir theilen daher hier eine Reihe betreffender Zeugnisse im Folgenden mit.

Martin Chemnitz schreibt in seinem „Examen Concilii Tridentini“: „Weil zum Kirchenamt viele Verrichtungen gehören, welche, wenn die Menge der Gläubigen sehr zahlreich ist, nicht wohl alle und jede von Einem oder Wenigen versehen werden können, so fing man in der apostolischen Zeit an, damit alles ordentlich, schicklich und zur Erbauung geschehe, wenn sich der Cötus der Kirche vervielfältigt hatte, jene Verrichtungen des Predigtamtes in gewisse Stufen\footnote{Aus dem folgenden geht hervor, daß Chemnitz nicht von Amtsstufen in episcopalistischem Sinne rede, daß er vielmehr nur ein Amt anerkennt, dessen verschiedene Functionen aber verschiedenen Personen als Gehilfen dessen, der das Amt κατ’ ἐξοχήν hat, übergeben und in diesem Sinne vermöge kirchlichen Ordnungsrechts zu Errichtung einer Art Amtsstufen gebraucht werden können.} von Kirchendienern einzutheilen, welche man τάξεις oder τάγματα hernach nannte, damit ein jeder seinen gewissen bestimmten Posten hätte, auf welchem er durch gewisse Verrichtungen des Predigtamts der Gemeinde diente. So besorgten im Anfange die Apostel das Amt des Wortes und der Sacramente, und zugleich auch die Austheilung und Verwaltung der Almosen. Hernach aber, als die Zahl der Jünger wuchs, übertrugen sie diesen Theil des Predigtamts, der die Almosen betrifft, andern, welche sie Diakonen d.\,i. Diener nannten. Und
\quellseite{24}
die Ursache, warum sie dies thun, geben sie an; damit sie nemlich am Amt des Worts und am Gebet ohne Abhaltungen anhalten könnten. Apost. 6, 4. Und dieser erste Ursprung der Stufen und Ordnungen des Predigtamts in der apostolischen Kirche zeigt, was für eine Ursache, was für eine Art und Weise, was für ein Zweck und Gebrauch sein solle von solchen Stufen und Ordnungen; „\gesperrt{damit nemlich nach Beschaffenheit der Kirchgemeine die einzelnen Verrichtungen, welche zum Predigtamt gehören, bequemer, richtiger, fleißiger und in Ordnung, mit einer gewissen Würde, zur Erbauung versehen werden}.“

Nachdem nun Chemnitz in dem Folgenden noch mehrere andere s.\,g. „Amtsstufen“ angeführt hat, als: Propheten, Evangelisten, Lehrer, Hirten, so fährt er weiter also fort: „1 Tim. 5, 17. thut Paulus zweier Gattungen von Presbytern Erwähnung, von denen die einen am Wort und an der Lehre arbeiteten, die andern den Kirchenzuchtsachen vorgesetzt waren; welches Presbyterats auch Tertullian im Apologeticus S. 39. gedenkt. Dieses sind ungefähr die Stufen, in welche, wie wir lesen, die Verrichtungen des Predigtamts zur Zeit der Apostel vertheilt waren. . . Es ist aber diese Erinnerung hinzuzusetzen: 1. daß es kein Gebot Gottes sei, welche oder wieviel solcher Stufen oder Ordnungen sein sollen. 2. Daß zu der Apostel Zeiten nicht in allen Gemeinen und immer dieselben und ebenso viele Stufen oder Ordnungen gewesen seien; was aus den Briefen Pauli, die an verschiedene Gemeinen geschrieben sind, zu schließen ist. 3. Es war zu der Apostel Zeit nicht eine solche Vertheilung jener Stufen, daß nicht öfters ein und derselbe alle diese Verrichtungen, welche zum Predigtamt gehören, übernahm und ausführte, was aus der apostolischen Geschichte bekannt ist. Solche Ordnungen waren daher zur Zeit der Apostel frei, indem man dabei auf Ordnung, Wohlständigkeit und Erbauung Rücksicht nahm; nur daß damals einige besondere Gaben, wie die der Sprachen, der Prophetie, des Apostolats und der Wunder gewissen Personen von Gott verliehen waren. Jene Stufen aber, von denen wir bisher geredet haben, waren nichts über und außer dem Amte des Wortes und der Sacramente, sondern eben die Verrichtungen des Predigtamts waren in jene Stufen eingetheilt, um der bereits aus einander gesetzten Ursachen willen. Jenes Beispiel der Apostel hat die erste Kirche auf dieselbe Art und Weise und mit gleicher Freiheit nachgeahmt. Man vertheilte die Grade der Verrichtungen des Predigtamtes; nicht aber auf ganz gleiche Weise, wie in der Corinthischen oder Ephessischen Gemeinde, sondern je nach den Umständen einer jeden Gemeinde. Woraus denn abzunehmen ist, was für eine Freiheit in Vertheilung jener Stufen stattfand. . . Auf diese Weise aber, zu diesem Endzweck und mit Beobachtung dieser Freiheit werden die meisten jener Stufen der alten Kirche auch bei uns beobachtet.“ (Part. II. Loc. 13. Sect. 2. de septem ordinibus, fol. 574. seqq.)

Johann Gerhard schreibt Folgendes: „In der apostolischen und
\quellseite{25}
ursprünglichen Kirche gab es zwei Gattungen von Presbytern, welche man lateinisch Seniores nennt, wie aus 1 Tim. 5, 17. geschlossen wird. Denn einige verwalteten das Lehramt, oder, wie der Apostel daselbst redet, arbeiteten im Wort und in der Lehre, welche Bischöfe, Pastoren etc. genannt wurden; andere aber waren nur für die Sittenzucht und Erhaltung der Kirchenzucht vorgesetzt, da die noch heidnische Obrigkeit die Lehrenden in der Kirche in diesem Stücke nicht unterstützte;\footnote{Aus dieser Bemerkung geht hervor, warum gerade in der lutherischen Kirche, wo die Obrigkeit lutherischen Bekenntnisses war, das Institut der Gemeindevorsteher nur ausnahmsweise sich findet. Zugleich liegt aber auch hierin ein Wink für uns hier in Amerika, wo die Obrigkeit unseres Bekenntnisses nicht ist, daß hier das Bedürfniß mitregierender Vorsteher um so mehr hervortritt.} diese wurden Regierer und Vorsteher genannt, wie aus 1 Cor. 12, 28. Röm. 12, 8. geschlossen wird. Ambrosius schreibt zu 1 Tim. 5. zu Anfang: „Auch die Synagoge und hernach die Kirche hat Seniores gehabt, ohne deren Rath nichts in der Kirche vorgenommen wurde, und ich weiß nicht, durch welche Nachlässigkeit dieses abgekommen ist, als etwa durch die Trägheit, oder vielmehr durch den Stolz der Lehrer, indem sie allein etwas gelten wollen.“ Beide Gattungen trugen gemeinschaftlich den Namen Vorsteher, 1 Tim. 5, 17., und Vorgesezte, Apost. 15, 22. Ebr. 13, 7. 17. 24. Aus beiden zugleich war jenes heilige Collegium gesammelt, welches Paulus das Presbyterium nennt, 1 Tim. 4, 14.: „Laß nicht aus der Acht die Gabe, die dir gegeben ist durch die Weissagung mit Handauflegung der Aeltesten“ (des Presbyteriums). Aus dieser Stelle wird geschlossen, daß bei der Ordination der Kirchendiener nicht allein die Pastoren, sondern auch die aus dem Volke gewählten Aeltesten im Namen der ganzen Kirche dem zu Ordinirenden die Hände aufgelegt haben,\footnote{Obwohl wir die historische Richtigkeit dieser Bemerkung Gerhards dahingestellt sein lassen, so machen wir doch hierbei darauf aufmerksam, daß ja freilich hiernach einem Gerhard ein ganz anderer Begriff von der Ordination vorgeschwebt haben müsse, als er jetzt wieder herrschend werden zu wollen scheint.} sowie 4 Mos. 8, 10. nicht allein Aaron, sondern auch die Aeltesten aus Israel bei der Ordination der Leviten denselben die Hand auflegten. (Einige jedoch meinen, daß hier unter dem Namen des Presbyteriums im speciellen Sinne die Versammlung derjenigen Aeltesten verstanden werde, welche Pastoren und Bischöfe waren.) Dem Presbyterium entspricht heut zu Tage das Consistorium oder der kirchliche Senat (Kirchenrath), in welchem kirchliche und politische Personen den Vorsitz haben, und vornehmlich für die Kirchenzucht Sorge tragen.“ (Loc. de ministerio § 232.) (Fortsetzung folgt.)
